\part*{Part I --- Signals, oscillators and damage}
\addcontentsline{toc}{part}{Part I --- Signals, oscillators and damage}

\section{Random signals}\label{sec:signals}

\subsection{Statistics of a signal}

Let $x_1,\dots,x_N$ be the samples of a signal. Its \emph{mean} and its \emph{central moments} of
order $m$ are
\begin{equation}
  \bar x = \frac1N\sum_{k=1}^N x_k, \qquad \mu_m = \frac1N\sum_{k=1}^N (x_k-\bar x)^m .
\end{equation}
The \emph{variance} is $\sigma^2 = \mu_2$, the standard deviation is $\sigma=\sqrt{\mu_2}$, and for a
zero-mean signal $\sigma$ is the root mean square (RMS) value. The \emph{skewness} $\mu_3/\sigma^3$ is
zero for a signal whose values are distributed symmetrically about the mean. The \emph{kurtosis}
(Pearson's definition, used everywhere in this manual and in the program) is
\begin{equation}\label{eq:kurtosis}
  \kappa = \frac{\mu_4}{\mu_2^{2}} .
\end{equation}
It measures the weight of the tails of the distribution of the values: rare large values raise
$\mu_4$ much more than $\mu_2$. The \emph{crest factor} is $\max_k|x_k-\bar x|/\sigma$.

The same definitions apply to a random variable $X$ when the averages over the samples are replaced
by expectations: $\mu_m = \E[(X-\E X)^m]$.

\subsection{The Gaussian distribution and its kurtosis}\label{sec:gausskurt}

A zero-mean Gaussian random variable $X$ with variance $\sigma^2$ has the probability density
\begin{equation}
  p(x) = \frac{1}{\sqrt{2\pi}\,\sigma}\, e^{-x^2/(2\sigma^2)} .
\end{equation}
Its even moments follow from one integration by parts. For an integer $m\ge 0$, write
$x^{m+2}e^{-x^2/(2\sigma^2)} = x^{m+1}\cdot\bigl(x\,e^{-x^2/(2\sigma^2)}\bigr)$ and note that
$x\,e^{-x^2/(2\sigma^2)}$ is the derivative of $-\sigma^2 e^{-x^2/(2\sigma^2)}$. Then
\begin{align}
  \E[X^{m+2}] &= \frac{1}{\sqrt{2\pi}\sigma}\int_{-\infty}^{\infty} x^{m+1}\, x\, e^{-x^2/(2\sigma^2)}\,\dd x
  \nonumber\\
  &= \frac{1}{\sqrt{2\pi}\sigma}\Bigl[-\sigma^2 x^{m+1} e^{-x^2/(2\sigma^2)}\Bigr]_{-\infty}^{\infty}
   + \frac{1}{\sqrt{2\pi}\sigma}\int_{-\infty}^{\infty} (m+1)x^{m}\,\sigma^2 e^{-x^2/(2\sigma^2)}\,\dd x
  \nonumber\\
  &= 0 + (m+1)\,\sigma^2\,\E[X^m] .
\end{align}
The bracket vanishes because the exponential decays faster than any power grows. Starting from
$\E[X^0]=1$: $\E[X^2] = 1\cdot\sigma^2\cdot 1 = \sigma^2$ and $\E[X^4] = 3\sigma^2\E[X^2] = 3\sigma^4$.
Therefore
\begin{result}[Kurtosis of a Gaussian variable]
\begin{equation}\label{eq:gauss4}
  \E[X^4] = 3\sigma^4, \qquad \kappa = \frac{3\sigma^4}{(\sigma^2)^2} = 3 .
\end{equation}
\end{result}
A signal with $\kappa>3$ is called \emph{leptokurtic}: it has more large values than a Gaussian
signal with the same RMS. Measured road and transport vibration is typically leptokurtic, with
$\kappa$ between 4 and 8, because its level changes with the road surface and contains shocks.

\subsection{Power spectral density}\label{sec:psd}

A record of duration $T$ sampled at $f_s$ can be written exactly as a finite Fourier series. With
$\Delta f = 1/T$ and the frequencies $f_m = m\,\Delta f$, $m = 0,1,\dots,N/2$,
\begin{equation}\label{eq:fourier}
  x(t_k) = \bar x + \sum_{m=1}^{N/2} A_m \cos(2\pi f_m t_k + \varphi_m),\qquad t_k = k\Delta t ,
\end{equation}
where $A_m\ge0$ is the amplitude and $\varphi_m$ the phase of the line $f_m$. A cosine of amplitude
$A_m$ has mean square value $A_m^2/2$ over a whole number of periods (because $\cos^2\theta =
\tfrac12 + \tfrac12\cos2\theta$ and the second term averages to zero), and two lines of different
frequencies are orthogonal over the record (the average of their product is zero). The variance of
the record is therefore the sum of the powers of its lines,
\begin{equation}\label{eq:parseval}
  \sigma^2 = \sum_{m} \frac{A_m^2}{2}.
\end{equation}
The \emph{one-sided power spectral density} (PSD) distributes this power over frequency: the power of
the line $f_m$ is spread over the band $\Delta f$ that the line represents,
\begin{equation}\label{eq:psddef}
  G(f_m) = \frac{A_m^2}{2\,\Delta f}, \qquad \sigma^2 = \sum_m G(f_m)\,\Delta f \;\approx\; \int_0^{f_s/2} G(f)\,\dd f .
\end{equation}
The unit of $G$ is (m/s$^2$)$^2$/Hz. For a random process the PSD is the expectation of this quantity
when $T\to\infty$; for a single record it must be \emph{estimated}, because the amplitudes of
individual lines scatter strongly from one record to another. The program uses Welch's estimate:
the record is cut into segments of length $1/r$ ($r$ is the requested resolution in Hz) overlapping
by 50\,\%, each segment is multiplied by a Hamming window (a raised-cosine taper that reduces the
leakage of power between lines), the squared moduli of their Fourier
transforms are scaled as in~\eqref{eq:psddef} and averaged. Averaging many segments reduces the
scatter; this is the only property of the estimate that the manual uses.

\subsection{Gaussian signals with a prescribed PSD}\label{sec:multisine}

Equation~\eqref{eq:fourier} also tells how to \emph{build} a signal with a prescribed PSD $G(f)$:
choose the amplitudes
\begin{equation}\label{eq:amplitudes}
  A_m = \sqrt{2\,G(f_m)\,\Delta f}
\end{equation}
and phases $\varphi_m$ drawn independently and uniformly in $[0,2\pi)$. This is a
\emph{random-phase multisine}. Its variance is $\sum_m A_m^2/2 = \sum_m G(f_m)\Delta f$, exactly the
area under the prescribed PSD, for every choice of the phases. The value $x(t)$ at a given time is a
sum of many independent terms $A_m\cos(2\pi f_m t+\varphi_m)$, each with zero mean and bounded; by the
central limit theorem the distribution of such a sum approaches a Gaussian distribution when many
lines carry comparable power. In practice a test of 60\,s has a line spacing of 1/60\,Hz, and
hundreds of lines fall in any band of interest, so the multisine is Gaussian for all purposes.

A second construction is exactly Gaussian for any number of lines. Write each line as
$A_m\cos(2\pi f_m t+\varphi_m) = U_m\cos(2\pi f_m t) - V_m\sin(2\pi f_m t)$ with $U_m =
A_m\cos\varphi_m$ and $V_m = A_m\sin\varphi_m$ (this is the addition formula of the cosine). If $U_m$
and $V_m$ are drawn as independent zero-mean Gaussian variables with variance $G(f_m)\Delta f$, then
$x(t)$ is, at every $t$, a sum of Gaussian variables with fixed coefficients, hence Gaussian, with
variance $\sum_m G(f_m)\Delta f\,(\cos^2+\sin^2) = \sum_m G(f_m)\Delta f$. The amplitude $A_m =
\sqrt{U_m^2+V_m^2}$ is then random instead of fixed: it is the length of a vector with two independent
Gaussian components, whose distribution (the Rayleigh distribution) is derived in
Section~\ref{sec:rayleigh},
with $\E[A_m^2/2] = G(f_m)\Delta f$. The pair $(U_m, V_m)$, written as the complex number $U_m+\ii V_m$,
is a \emph{complex normal coefficient}; the program draws such coefficients with unit mean square
modulus and scales them to the power of each line.

Keeping the phases (or the coefficients) of a realisation fixed while changing the amplitudes
$A_m$ is called \emph{frozen phases}: the signal then changes smoothly with its PSD, which is what
the closed-loop corrections of Section~\ref{sec:closedloop} rely on.

\subsection{Modulated Gaussian signals}\label{sec:modulated}

A leptokurtic signal can be built by multiplying a stationary Gaussian \emph{carrier} $g(t)$ (zero
mean, variance $\sigma_g^2$) by a positive \emph{envelope} $e(t)$ that varies slowly compared with the
carrier:
\begin{equation}\label{eq:envcarrier}
  x(t) = e(t)\,g(t).
\end{equation}
Divide the duration $T$ into windows short enough for $e$ to be practically constant in each
(window $w$ has duration $T_w$ and envelope value $e_w$) but long enough to contain many periods of
the carrier, so that the time average of $g^2$ and $g^4$ over a window equals its expectation.
Using~\eqref{eq:gauss4} for the carrier,
\begin{align}
  \tavg{x^2} &= \sum_w \frac{T_w}{T}\, e_w^2\, \sigma_g^2 = \sigma_g^2\,\tavg{e^2}, \\
  \tavg{x^4} &= \sum_w \frac{T_w}{T}\, e_w^4\, 3\sigma_g^4 = 3\sigma_g^4\,\tavg{e^4},
\end{align}
and the kurtosis~\eqref{eq:kurtosis} of the product is
\begin{result}[Kurtosis of a modulated Gaussian signal]
\begin{equation}\label{eq:kmod}
  \kappa = \frac{3\sigma_g^4\tavg{e^4}}{\bigl(\sigma_g^2\tavg{e^2}\bigr)^2} = 3\,\frac{\tavg{e^4}}{\tavg{e^2}^2}
  = 3\,\frac{\tavg{s^2}}{\tavg{s}^2}, \qquad s(t) = e(t)^2\sigma_g^2 .
\end{equation}
\end{result}
The quantity $s(t)$ is the \emph{local variance} of the signal. The last form of~\eqref{eq:kmod}
states a general rule used throughout Part~III: if, given a slowly varying local variance $s(t)$,
the signal is Gaussian, its kurtosis is three times the ratio of the mean square of $s$ to the square
of its mean. By the Cauchy--Schwarz inequality, $\tavg{s}^2 \le \tavg{s^2}\tavg{1}$, so
$\kappa\ge3$, with equality only for a constant $s$. Modulation can raise the kurtosis but never lower
it below the Gaussian value.
